\frfilename{mt256.tex}
\versiondate{6.8.15}
\copyrightdate{2000}

\def\chaptername{Ukuran produk}
\def\sectionname{Ukuran Radon pada $\BbbR^r$}

\newsection{256}

\def\headlinesectionname{Ukuran Radon pada {$\eightBbb R^r$}}

Pada bagian berikutnya, dan kembali dalam Bab 27 dan 28,
kita perlu membahas kelas utama ukuran pada ruang Euklides.
Untuk pembahasan yang semestinya mengenai kelas ini, dan keterkaitan
antara ukuran-ukuran dengan topologi-topologi yang terlibat, kita harus menunggu sampai
Jilid 4.   Oleh karena itu, untuk saat ini saya menyajikan definisi-definisi yang disesuaikan dengan
kasus yang sedang dihadapi, dengan memperingatkan Anda bahwa perumuman yang tepat tidak
sepenuhnya jelas.   Saya memberikan definisi (256A) dan suatu karakterisasi
(256C) untuk ukuran Radon pada ruang Euklides, serta teorema-teorema tentang
konstruksi
ukuran Radon sebagai integral tak tentu (256E, 256J), sebagai ukuran citra
(256G), dan sebagai ukuran produk (256K).   Sambil lalu saya memberikan suatu versi
teorema Lusin mengenai fungsi terukur pada ruang ukur Radon
(256F).

Di seluruh bagian ini, $r$ dan $s$ adalah bilangan bulat yang lebih besar daripada atau sama
dengan $1$.

\vleader{48pt}{256A}{Definisi} Misalkan $\nu$ suatu ukuran pada $\BbbR^r$
dan $\Sigma$ domainnya.

\spheader 256Aa $\nu$ adalah suatu {\bf ukuran topologis} jika setiap himpunan terbuka
termasuk dalam $\Sigma$.   \cmmnt{Perhatikan bahwa dalam hal ini setiap himpunan Borel,
dan khususnya setiap himpunan tertutup, termasuk dalam $\Sigma$.}

\spheader 256Ab $\nu$ adalah {\bf berhingga secara lokal} jika setiap himpunan terbatas mempunyai
ukuran luar berhingga.

\spheader 256Ac Jika $\nu$ adalah ukuran topologis, ukuran itu {\bf regular
dalam terhadap himpunan kompak} jika

\Centerline{$\nu E=\sup\{\nu K:K\subseteq E$ kompak$\}$}

\noindent untuk setiap $E\in\Sigma$.   \cmmnt{(Karena $\nu$ adalah
ukuran topologis, dan himpunan kompak itu tertutup (2A2Ec), $\nu K$
terdefinisi untuk setiap himpunan kompak $K$.)}

\spheader 256Ad $\nu$ adalah suatu {\bf ukuran Radon} jika ukuran itu merupakan ukuran topologis
lengkap yang berhingga secara lokal dan regular dalam terhadap
himpunan kompak.


\leader{256B}{}\cmmnt{ Kita akan terbantu jika dapat menggunakan
fakta-fakta elementer berikut.

\medskip

\noindent}{\bf Lema} Misalkan $\nu$ suatu ukuran Radon pada $\BbbR^r$, dan
$\Sigma$ domainnya.

(a) $\nu$ adalah $\sigma$-hingga.

(b) Untuk setiap $E\in\Sigma$ dan setiap
$\epsilon>0$ terdapat suatu himpunan tertutup
$F\subseteq E$ dan suatu himpunan terbuka $G\supseteq E$ sedemikian sehingga
$\nu(G\setminus F)\le\epsilon$.

(c) Untuk setiap $E\in\Sigma$ terdapat suatu himpunan
$H\subseteq E$, yang dapat dinyatakan sebagai gabungan suatu barisan himpunan kompak,
sedemikian sehingga $\nu(E\setminus H)=0$.

(d) Setiap fungsi kontinu bernilai real pada $\BbbR^r$ terukur
terhadap $\Sigma$.

(e) Jika $h:\BbbR^r\to\Bbb R$ kontinu dan mempunyai dukungan terbatas,
maka $h$ terintegralkan terhadap $\nu$.

\proof{{\bf (a)} Untuk setiap $n\in\Bbb N$, $B(\tbf{0},n)=\{x:\|x\|\le n\}$
adalah suatu himpunan tertutup dan terbatas,
sehingga merupakan himpunan Borel.   Jadi jika $\nu$ adalah suatu ukuran Radon pada $\BbbR^r$,
$\sequencen{B(\tbf{0},n)}$ adalah suatu barisan
himpunan dengan ukuran berhingga yang menutupi $\BbbR^r$.

\medskip

{\bf (b)} Ambil $E_n=\{x:x\in E,\,n\le\|x\|<n+1\}$ untuk setiap $n$.   Maka
$\nu E_n<\infty$, sehingga terdapat suatu himpunan kompak $K_n\subseteq E_n$ sedemikian sehingga
$\nu K_n\ge\nu E_n-2^{-n-2}\epsilon$.   Ambil $F=\bigcup_{n\in\Bbb N}K_n$;
maka

\Centerline{$\nu(E\setminus F)
=\sum_{n=0}^{\infty}\nu(E_n\setminus K_n)\le\Bover12\epsilon$.}

\noindent Selain itu $F\subseteq E$ dan $F$ tertutup karena

\Centerline{$F\cap B(\tbf{0},n)=\bigcup_{i\le n}K_i\cap B(\tbf{0},n)$}

\noindent tertutup untuk setiap $n$.

Dengan cara yang sama, terdapat suatu himpunan tertutup $F'\subseteq\BbbR^r\setminus E$
sedemikian sehingga $\nu((\BbbR^r\setminus E)\setminus F')\le\bover12\epsilon$.
Dengan mengambil $G=\BbbR^r\setminus F'$, kita melihat bahwa $G$ terbuka, bahwa
$G\supseteq E$, dan bahwa $\nu(G\setminus E)\le\bover12\epsilon$, sehingga
$\nu(G\setminus F)\le\epsilon$, sebagaimana dikehendaki.

\medskip

{\bf (c)} Berdasarkan (b), untuk setiap $n\in\Bbb N$ kita dapat memilih suatu himpunan tertutup
$F_n\subseteq E$ sedemikian sehingga $\nu(E\setminus F_n)\le 2^{-n}$.   Ambil
$H=\bigcup_{n\in\Bbb N}F_n$;  maka $H\subseteq E$ dan $\nu(E\setminus
H)=0$, dan juga $H=\bigcup_{m,n\in\Bbb N}B(\tbf{0},m)\cap F_n$ merupakan suatu
gabungan terhitung dari himpunan-himpunan kompak.

\medskip

{\bf (d)} Jika $h:\BbbR^r\to\Bbb R$ kontinu, semua himpunan
$\{x:h(x)>a\}$ terbuka, sehingga termasuk dalam $\Sigma$.

\medskip

{\bf (e)} Berdasarkan (d), $h$ terukur.   Sekarang kita mengandaikan bahwa terdapat
suatu $n\in\Bbb N$ sedemikian sehingga $h(x)=0$ kapan pun $x\notin B(\tbf{0},n)$.
Karena $B(\tbf{0},n)$ kompak (2A2F), $h$ terbatas pada $B(\tbf{0},n)$
(2A2G), dan kita mempunyai $|h|\le \gamma\chi B(\tbf{0},n)$ untuk suatu $\gamma$;
karena $\nu B(\tbf{0},n)$ berhingga, $h$ terintegralkan terhadap $\nu$.
}%end of proof of 256B

\leader{256C}{Teorema} Suatu ukuran $\nu$ pada $\BbbR^r$ adalah ukuran Radon
jika dan hanya jika ukuran itu merupakan pelengkapan suatu ukuran berhingga secara lokal yang terdefinisi pada
aljabar-sigma\cmmnt{ $\Cal B$} dari subhimpunan-subhimpunan Borel $\BbbR^r$.

\proof{{\bf (a)} Misalkan terlebih dahulu bahwa $\nu$ adalah suatu ukuran Radon.   Tuliskan
$\Sigma$ untuk domainnya.

\medskip

\quad{\bf (i)} Ambil $\nu_0=\nu\restr\Cal B$.   Maka $\nu_0$ adalah suatu ukuran
dengan domain $\Cal B$, dan ukuran itu berhingga secara lokal karena
$\nu_0B(\tbf{0},n)=\nu B(\tbf{0},n)$ berhingga untuk setiap $n$.   Misalkan
$\hat\nu_0$ pelengkapan dari $\nu_0$ (212C).

\medskip

\quad{\bf (ii)} Jika menurut $\hat\nu_0$ himpunan $E$ terukur, terdapat $E_1$,
$E_2\in\Cal B$ sedemikian sehingga $E_1\subseteq E\subseteq E_2$ dan
$\nu_0(E_2\setminus E_1)=0$.   Sekarang $E\setminus E_1\subseteq E_2\setminus
E_1$ harus
terabaikan terhadap $\nu$;  karena $\nu$ lengkap, $E\in\Sigma$ dan

\Centerline{$\nu E=\nu E_1=\nu_0 E_1=\hat\nu_0E$.}

\medskip

\quad{\bf (iii)} Jika $E\in\Sigma$, maka berdasarkan 256Bc terdapat suatu
himpunan Borel $H\subseteq E$ sedemikian sehingga $\nu(E\setminus H)=0$.   Demikian pula,
terdapat suatu himpunan Borel $H'\subseteq\BbbR^r\setminus E$ sedemikian sehingga
$\nu((\BbbR^r\setminus E)\setminus H')=0$, sehingga kita mempunyai $H\subseteq
E\subseteq\BbbR^r\setminus H'$ dan

\Centerline{$\nu_0((\BbbR^r\setminus H')\setminus H)=\nu((\Bbb
R^r\setminus H')\setminus H)=0$.}

\noindent Jadi $\hat\nu_0E$ terdefinisi dan sama dengan $\nu_0H=\nu E$.

Ini menunjukkan bahwa $\nu=\hat\nu_0$ adalah pelengkapan dari ukuran Borel berhingga secara lokal
$\nu\restr\Cal B$.   Dan ini berlaku untuk setiap ukuran Radon
$\nu$ pada $\BbbR^r$.

\medskip

{\bf (b)} Untuk bagian selebihnya dari pembuktian ini, saya mengandaikan bahwa $\nu_0$ adalah suatu ukuran berhingga
secara lokal dengan domain $\Cal B$ dan $\nu$ adalah pelengkapannya.    Tuliskan
$\Sigma$ untuk domain $\nu$.   Kita menyebut suatu subhimpunan $\BbbR^r$
sebagai suatu {\bf himpunan K$_{\sigma}$} jika subhimpunan itu dapat dinyatakan sebagai gabungan suatu
barisan himpunan kompak.   Perhatikan bahwa setiap himpunan K$_{\sigma}$ adalah himpunan Borel,
sehingga termasuk dalam $\Sigma$.   Ambil

\Centerline{$\Cal A=\{E:E\in\Sigma$, terdapat suatu himpunan K$_{\sigma}$
$H\subseteq E$ sedemikian sehingga $\nu(E\setminus H)=0\}$,}

\Centerline{$\Cal C=\{E:E\in\Cal A,\,\BbbR^r\setminus E\in\Cal A\}$.}

\medskip

{\bf (c)(i)} Setiap himpunan terbuka itu sendiri merupakan suatu himpunan K$_{\sigma}$, sehingga termasuk dalam
$\Cal A$.
\Prf\ Misalkan $G\subseteq\BbbR^r$ terbuka.   Jika $G=\emptyset$ maka $G$
kompak dan hasilnya langsung.   Jika tidak, misalkan $\Cal I$ himpunan
interval tertutup berbentuk $[q,q']$, dengan $q$, $q'\in\Bbb Q^r$,
yang termuat dalam $G$.   Maka semua anggota $\Cal I$ tertutup
dan terbatas, sehingga kompak.   Jika $x\in G$, terdapat suatu $\delta>0$
sedemikian sehingga $B(x,\delta)=\{y:\|y-x\|\le\delta\}\subseteq G$;  sekarang terdapat
suatu $I\in\Cal I$ sedemikian sehingga $x\in I\subseteq B(x,\delta)$.   Dengan demikian
$G=\bigcup\Cal I$.   Tetapi $\Cal I$ terhitung, sehingga $G$ adalah K$_{\sigma}$.\
\Qed

\medskip

\quad{\bf (ii)} Setiap subhimpunan tertutup dari $\Bbb R^r$ adalah K$_{\sigma}$, sehingga
termasuk dalam $\Cal A$.
\Prf\ Jika $F\subseteq\Bbb R^r$ tertutup, maka
$F=\bigcup_{n\in\Bbb N}F\cap B(\tbf{0},n)$;  tetapi setiap $F\cap
B(\tbf{0},n)$ tertutup dan terbatas, sehingga kompak.\ \Qed

\medskip

\quad{\bf (iii)} Jika $\sequencen{E_n}$ adalah sebarang barisan dalam $\Cal A$, maka
$E=\bigcup_{n\in\Bbb N}E_n$ termasuk dalam
$\Cal A$.   \Prf\ Untuk setiap $n\in\Bbb N$ kita mempunyai suatu keluarga terhitung
$\Cal K_n$ yang terdiri atas subhimpunan kompak dari $E_n$ sedemikian sehingga $\nu(E_n\setminus\bigcup\Cal K_n)=0$;  sekarang
$\Cal K=\bigcup_{n\in\Bbb N}\Cal K_n$ adalah suatu keluarga terhitung
yang terdiri atas subhimpunan kompak dari $E$, dan
$E\setminus\bigcup\Cal K
\subseteq\bigcup_{n\in\Bbb N}(E_n\setminus\bigcup\Cal K_n)$ bersifat
terabaikan terhadap $\nu$.\ \Qed

\medskip

\quad{\bf (iv)} Jika $\sequencen{E_n}$ adalah sebarang barisan dalam $\Cal A$, maka
$F=\bigcap_{n\in\Bbb N}E_n\in\Cal A$.   \Prf\ Untuk setiap $n\in\Bbb N$, misalkan
$\sequence{i}{K_{ni}}$ suatu barisan yang terdiri atas subhimpunan kompak dari $E_n$ sedemikian
sehingga $\nu(E_n\setminus\bigcup_{i\in\Bbb N}K_{ni})=0$.   Ambil
$K'_{nj}=\bigcup_{i\le j}K_{ni}$ untuk setiap $j$, sehingga

\Centerline{$\nu(E_n\cap H)=\lim_{j\to\infty}\nu(K'_{nj}\cap H)$}

\noindent untuk setiap $H\in\Sigma$.   Sekarang, untuk setiap $m$, $n\in\Bbb N$,
pilih $j(m,n)$ sedemikian sehingga

\Centerline{$\nu(E_n\cap B(\tbf{0},m)\cap K'_{n,j(m,n)})
\ge\nu(E_n\cap B(\tbf{0},m))-2^{-(m+n)}$.}

\noindent Ambil $K_m=\bigcap_{n\in\Bbb N}K'_{n,j(m,n)}$;  maka $K_m$
tertutup (karena merupakan irisan himpunan-himpunan tertutup) dan terbatas (karena merupakan
subhimpunan $K'_{0,j(m,0)}$), sehingga kompak.   Selain itu $K_m\subseteq F$,
karena $K'_{n,j(m,n)}\subseteq E_n$ untuk setiap $n$, dan

\Centerline{$\nu(F\cap B(\tbf{0},m)\setminus K_m)
\le\sum_{n=0}^{\infty}\nu(E_n\cap B(\tbf{0},m)\setminus K'_{n,j(m,n)})
\le\sum_{n=0}^{\infty}2^{-(m+n)}
=2^{-m+1}$.}

\noindent Akibatnya $H=\bigcup_{m\in\Bbb N}K_m$ adalah suatu subhimpunan K$_{\sigma}$
dari $F$ dan

\Centerline{$\nu(F\cap B(\tbf{0},m)\setminus H)
\le\inf_{k\ge m}\nu(F\cap B(\tbf{0},k)\setminus K_k)
=0$}

\noindent untuk setiap $m$, sehingga $\nu(F\setminus H)=0$ dan $F\in\Cal A$.\ \Qed

\medskip

{\bf (d)} $\Cal C$ adalah suatu aljabar-sigma dari subhimpunan $\Bbb R^r$.
\Prf\ {(i)} $\emptyset$ dan komplemennya terbuka, sehingga termasuk dalam $\Cal
A$ dan dengan demikian termasuk dalam $\Cal C$.
{(ii)} Jika $E\in\Cal C$ maka $\BbbR^r\setminus E$ dan $\Bbb
R^r\setminus(\BbbR^r\setminus E)=E$ termasuk dalam $\Cal A$, sehingga $\Bbb
R^r\setminus E\in\Cal C$.   {(iii)} Misalkan $\sequencen{E_n}$ suatu
barisan dalam $\Cal C$ dengan gabungan $E$.   Berdasarkan (c-iii) dan (c-iv),

\Centerline{$E\in\Cal A$,
\quad$\BbbR^r\setminus E
=\bigcap_{n\in\Bbb N}(\BbbR^r\setminus E_n)\in\Cal A$,}

\noindent sehingga $E\in\Cal C$.\  \Qed

\medskip

{\bf (e)} Berdasarkan (c-i) dan (c-ii), setiap himpunan terbuka termasuk dalam $\Cal C$;
akibatnya setiap himpunan Borel
termasuk dalam $\Cal C$ dan dengan demikian termasuk dalam $\Cal A$.   Sekarang jika $E$ adalah sebarang
anggota
dari $\Sigma$, terdapat suatu himpunan Borel $E_1\subseteq E$ sedemikian sehingga
$\nu(E\setminus E_1)=0$ dan suatu himpunan K$_{\sigma}$ $H\subseteq E_1$ sedemikian
sehingga $\nu(E_1\setminus H)=0$.   Nyatakan $H$ sebagai $\bigcup_{n\in\Bbb N}K_n$
dengan setiap $K_n$ kompak;  maka

\Centerline{$\nu E
=\nu H
=\lim_{n\to\infty}\nu(\bigcup_{i\le n}K_i)
\le\sup_{K\subseteq E\text{ kompak}}\nu K
\le\nu E$}

\noindent karena $\bigcup_{i\in n}K_i$ adalah suatu subhimpunan kompak dari $E$ untuk
setiap $n$.

\medskip

{\bf (f)} Dengan demikian $\nu$ regular dalam terhadap himpunan kompak.
Namun tentu saja ukuran itu lengkap (karena merupakan pelengkapan $\nu_0$) dan merupakan ukuran topologis
yang berhingga secara
lokal (karena $\nu_0$ demikian);  jadi ukuran itu merupakan suatu ukuran Radon.
Ini menyelesaikan pembuktian.
}%end of proof of 256C

\leader{256D}{Proposisi} Jika $\nu$ dan $\nuprime$ adalah dua ukuran Radon
pada $\BbbR^r$, pernyataan-pernyataan berikut berekuivalen:

(i) $\nu=\nuprime$;

(ii) $\nu K=\nuprime K$ untuk setiap himpunan kompak $K\subseteq\BbbR^r$;

(iii) $\nu G=\nuprime G$ untuk setiap himpunan terbuka $G\subseteq\BbbR^r$;

(iv) $\int h\,d\nu=\int h\,d\nuprime$ untuk setiap fungsi kontinu
$h:\BbbR^r\to\Bbb R$ yang mempunyai dukungan terbatas.

\proof{{\bf (a)(i)$\Rightarrow$(iv)} langsung.

\medskip

{\bf (b)(iv)$\Rightarrow$(iii)} Jika (iv) benar, dan $G\subseteq\BbbR^r$
adalah suatu himpunan terbuka, maka untuk setiap $n\in\Bbb N$ ambil

\Centerline{$h_n(x)
=\min(1,2^n\inf_{y\in\BbbR^r\setminus(G\cap B(\tbf{0},n))}\|y-x\|)$}

\noindent untuk $x\in\BbbR^r$.   Maka $h_n$ kontinu (bahkan
$|h_n(x)-h_n(x')|\le 2^n\|x-x'\|$ untuk semua $x$, $x'\in\BbbR^r$) dan bernilai nol
di luar $B(\tbf{0},n)$, sehingga $\int h_nd\nu=\int h_nd\nuprime$.   Selanjutnya,
$\sequencen{h_n(x)}$ adalah suatu barisan tak menurun yang konvergen ke
$\chi G(x)$ untuk setiap $x\in\BbbR^r$.   Jadi

\Centerline{$\nu G=\lim_{n\to\infty}\int h_nd\nu
=\lim_{n\to\infty}\int h_nd\nuprime=\nuprime G$,}

\noindent berdasarkan 135Ga.   Karena $G$ sebarang, (iii) benar.

\medskip

{\bf (c)(iii)$\Rightarrow$(ii)} Jika (iii) benar, dan
$K\subseteq\BbbR^r$ kompak, pilih $n$ cukup besar sehingga $\|x\|<n$ untuk
setiap $x\in K$.   Ambil $G=\{x:\|x\|<n\}$, $H=G\setminus K$.   Maka $G$
dan $H$ terbuka dan $G$ terbatas, sehingga $\nu G=\nuprime G$ berhingga,
dan

\Centerline{$\nu K=\nu G-\nu H=\nuprime G-\nuprime H=\nuprime K$.}

\noindent Karena $K$ sebarang, (ii) benar.

\medskip

{\bf (d)(ii)$\Rightarrow$(i)} Jika $\nu$ dan $\nuprime$ bernilai sama pada himpunan-himpunan
kompak, maka

\Centerline{$\nu E=\sup_{K\subseteq E\text{ kompak}}\nu K
=\sup_{K\subseteq E\text{ kompak}}\nuprime K
=\nuprime E$}

\noindent untuk setiap himpunan Borel $E$.   Jadi
$\nu\restr\Cal B=\nuprime\restr\Cal B$, dengan $\Cal B$ aljabar himpunan-himpunan Borel.
Tetapi karena $\nu$ dan $\nuprime$ keduanya merupakan pelengkapan dari
pembatasan masing-masing pada $\Cal B$, keduanya identik.
}%end of proof of 256D

\leader{256E}{}\cmmnt{ Kiranya sudah waktunya saya memberikan beberapa contoh
ukuran Radon.   Namun, beberapa baris dapat dihemat jika terlebih dahulu saya menetapkan
sejumlah konstruksi dasar.   Pada tahap ini Anda mungkin ingin melihat sekilas
256H.

\medskip

\noindent}{\bf Teorema} Misalkan $\nu$ suatu ukuran Radon pada $\BbbR^r$, dengan
domain $\Sigma$, dan
$f$ suatu fungsi nonnegatif yang terukur terhadap $\Sigma$, yang didefinisikan pada suatu
subhimpunan yang koterabaikan terhadap $\nu$ di dalam $\BbbR^r$.   Andaikan bahwa $f$
{\bf terintegralkan secara lokal} dalam arti bahwa $\int_Efd\nu<\infty$ untuk setiap
himpunan terbatas $E$.
Maka ukuran integral tak tentu $\nuprime$ pada $\BbbR^r$ yang didefinisikan dengan
syarat bahwa

\Centerline{$\nuprime E=\int_Efd\nu$ kapan pun
$E\cap\{x:x\in\dom f,\,f(x)>0\}\in\Sigma$}

\noindent merupakan ukuran Radon pada $\BbbR^r$.

\proof{ Untuk konstruksi $\nuprime$, lihat 234I-234L.
Ukuran integral tak tentu, sebagaimana saya definisikan, selalu lengkap (234I).
$\nuprime$ berhingga secara lokal karena $f$
terintegralkan secara lokal.   $\nuprime$ merupakan
ukuran topologis karena setiap himpunan terbuka termasuk dalam $\Sigma$ dan
oleh karena itu termasuk dalam domain $\Sigma'$ dari $\nuprime$.
Untuk melihat bahwa $\nuprime$ regular dalam terhadap himpunan
kompak, ambil sebarang himpunan $E\in\Sigma'$, dan tetapkan
$E'=\{x:x\in E\cap\dom f,\,f(x)>0\}$.
Maka $E'\in\Sigma$, sehingga terdapat himpunan $H\subseteq E'$,
yang dapat dinyatakan sebagai gabungan suatu barisan himpunan kompak,
sedemikian sehingga $\nu(E'\setminus H)=0$.   Dalam hal ini

\Centerline{$\nuprime(E\setminus H)=\int_{E\setminus H}fd\nu=0$.}

\noindent Misalkan $\sequencen{K_n}$ suatu barisan himpunan kompak dengan gabungan
$H$;  maka

\Centerline{$\nuprime E=\nuprime H
=\lim_{n\to\infty}\nuprime(\bigcup_{i\le n}K_i)
\le\sup_{K\subseteq E\text{ kompak}}\nuprime K\le\nuprime E$.}

\noindent Karena $E$ sebarang, $\nuprime$ regular dalam terhadap
himpunan kompak.
}%end of proof of 256E

\leader{256F}{Teorema} Misalkan $\nu$ suatu ukuran Radon pada $\BbbR^r$, dan
$\Sigma$ domainnya.   Misalkan $f:D\to\Bbb R$ suatu fungsi yang terukur terhadap $\Sigma$,
dengan $D\subseteq\BbbR^r$.   Maka untuk setiap $\epsilon>0$
terdapat himpunan tertutup $F\subseteq\BbbR^r$ sedemikian sehingga
$\nu(\BbbR^r\setminus F)\le\epsilon$ dan $f\restr F$ kontinu.

\proof{ Menurut 121I, terdapat suatu fungsi yang terukur terhadap $\Sigma$,
yakni $h:\BbbR^r\to\Bbb R$, yang memperluas $f$.
Tuliskan unsur-unsur $\Bbb Q$ sebagai barisan $\sequencen{q_n}$.
Untuk setiap $n\in\Bbb N$ tetapkan $E_n=\{x:h(x)\le q_n\}$, $E'_n=\{x:h(x)>q_n\}$
dan gunakan 256Bb untuk memilih himpunan-himpunan tertutup $F_n\subseteq E_n$,
$F'_n\subseteq E'_n$ sedemikian sehingga
$\nu(E_n\setminus F_n)\le 2^{-n-2}\epsilon$ dan
$\nu(E'_n\setminus F'_n)\le 2^{-n-2}\epsilon$.   Tetapkan
$F=\bigcap_{n\in\Bbb N}(F_n\cup F'_n)$;  maka $F$ tertutup dan

\Centerline{$\nu(\BbbR^r\setminus F)
\le\sum_{n=0}^{\infty}\nu(\BbbR^r\setminus(F_n\cup F'_n))
\le\sum_{n=0}^{\infty}\nu(E_n\setminus F_n)+\nu(E'_n\setminus F'_n)
\le\epsilon$.}

Saya akan menunjukkan bahwa $h\restr F$ kontinu.   \Prf\ Andaikan $x\in F$
dan $\delta>0$.   Maka terdapat $m$, $n\in\Bbb N$ sedemikian sehingga

\Centerline{$h(x)-\delta\le q_m<h(x)\le q_n\le h(x)+\delta$.}

\noindent Ini berarti bahwa $x\in E'_m\cap E_n$;  akibatnya
$x\notin F_m\cup F'_n$.
Karena $F_m\cup F'_n$ tertutup, terdapat $\eta>0$
sedemikian sehingga $y\notin F_m\cup F'_n$ kapan pun $\|y-x\|\le\eta$.
Sekarang andaikan bahwa $y\in F$ dan $\|y-x\|\le\eta$.   Maka
$y\in(F_m\cup F'_m)\cap(F_n\cup F'_n)$ dan $y\notin F_m\cup F'_n$,
sehingga $y\in F'_m\cap F_n\subseteq E'_m\cap E_n$ dan $q_m<h(y)\le q_n$.
Akibatnya
$|h(y)-h(x)|\le\delta$.   Karena $x$ dan $\delta$ sebarang, $h\restr F$
kontinu.\ \QeD\
Dengan demikian $f\restr F=(h\restr F)\restr D$ kontinu, sebagaimana diperlukan.
}%end of proof of 256F

\leader{256G}{Teorema} Misalkan $\nu$ suatu ukuran Radon pada $\BbbR^r$, dengan
domain $\Sigma$, dan
andaikan bahwa $\phi:\BbbR^r\to\BbbR^s$ terukur dalam arti bahwa
semua koordinatnya terukur terhadap $\Sigma$.   Jika ukuran citra
$\nuprime=\nu\phi^{-1}$\cmmnt{ (234D)} berhingga secara lokal, maka ukuran itu merupakan
ukuran Radon.

\proof{ Tuliskan $\Sigma'$ untuk
domain $\nuprime$.   Jika $\phi=(\phi_1,\ldots,\phi_s)$, maka

\Centerline{$\phi^{-1}[\{y:\eta_j\le\alpha\}]
=\{x:\phi_j(x)\le\alpha\}\in\Sigma$,}

\noindent sehingga $\{y:\eta_j\le\alpha\}\in\Sigma'$
untuk setiap $j\le s$, $\alpha\in\Bbb R$; di sini saya menulis
$y=(\eta_1,\ldots,\eta_s)$ untuk $y\in\BbbR^s$.   Akibatnya setiap
subhimpunan Borel dari $\BbbR^s$ termasuk dalam $\Sigma'$ (121J), dan $\nuprime$
merupakan ukuran topologis.   Ukuran itu lengkap menurut 234Eb.

Hal yang penting tentu saja ialah bahwa $\nuprime$ regular dalam terhadap
himpunan
kompak.   \Prf\ Andaikan bahwa $F\in\Sigma'$ dan bahwa $\gamma<\nuprime F$.
Untuk setiap $j\le s$, terdapat himpunan tertutup $H_j\subseteq\BbbR^r$ sedemikian
sehingga $\phi_j\restr H_j$ kontinu dan
$\nu(\BbbR^r\setminus H_j)<\bover1{s}(\nuprime F-\gamma)$, menurut 256F.   Tetapkan $H=\bigcap_{j\le s}H_j$;  maka $H$ tertutup dan $\phi\restr H$ kontinu dan

\Centerline{$\nu(\BbbR^r\setminus H)<\nuprime F-\gamma
=\nu\phi^{-1}[F]-\gamma$,}

\noindent sehingga
$\nu(\phi^{-1}[F]\cap H)>\gamma$.   Ambil himpunan kompak $K\subseteq\phi^{-1}[F]\cap H$
sedemikian sehingga $\nu K\ge\gamma$, dan tetapkan $L=\phi[K]$.
Karena $K\subseteq H$ dan $\phi\restr H$ kontinu, $L$ kompak
(2A2Eb).   Tentu saja $L\subseteq F$, dan

\Centerline{$\nuprime L=\nu\phi^{-1}[L]\ge\nu K\ge\gamma$.}

\noindent Karena $F$ dan $\gamma$ sebarang, $\nuprime$ regular dalam
terhadap himpunan kompak.\ \Qed

Karena $\nuprime$ berhingga secara lokal menurut hipotesis teorema, ukuran itu
merupakan ukuran Radon.
}%end of proof of 256G

\leader{256H}{Contoh-contoh}\cmmnt{ Akhirnya saya sampai pada contoh-contoh yang
dijanjikan.

\medskip

} {\bf (a)} Ukuran Lebesgue pada $\BbbR^r$ merupakan ukuran Radon.
\prooflet{(Ukuran itu merupakan ukuran topologis menurut 115G, dan regular dalam terhadap
himpunan kompak menurut 134Fb.)}

\spheader 256Hb Suatu ukuran yang bertumpu pada titik di $\BbbR^r$ merupakan ukuran Radon
jika dan hanya jika ukuran itu berhingga secara lokal.   \prooflet{\Prf\ Misalkan $\mu$ suatu ukuran yang bertumpu pada titik
pada $\BbbR^r$.   Jika ukuran itu merupakan ukuran Radon, tentu saja ukuran itu
berhingga secara lokal.   Jika ukuran itu berhingga secara lokal, maka ukuran itu jelas merupakan
ukuran topologis
lengkap, karena ukuran itu mengukur setiap subhimpunan $\BbbR^r$.   Misalkan
$h:\BbbR^r\to[0,\infty]$ sedemikian sehingga $\mu E=\sum_{x\in E}h(x)$ untuk setiap
$E\subseteq\BbbR^r$.   Ambil sebarang $E\subseteq\BbbR^r$.   Maka

$$\eqalign{\mu E
&=\sum_{x\in E}h(x)
=\sup_{I\subseteq E\text{ berhingga}}\sum_{x\in I}h(x)\cr&
=\sup_{I\subseteq E\text{ berhingga}}\mu I
\le\sup_{K\subseteq E\text{ kompak}}\mu K
\le\mu E
\cr}$$

\noindent sehingga $\mu E=\sup_{K\subseteq E\text{ kompak}}\mu K$;
dengan demikian $\mu$ regular dalam terhadap himpunan kompak dan merupakan
ukuran Radon.\ \Qed}%end of prooflet

\spheader 256Hc \cmmnt{Sekarang kita sampai pada suatu gagasan baru.}  Ingat bahwa
himpunan Cantor $C$\cmmnt{ (134G)} merupakan subhimpunan tertutup dari $[0,1]$ yang terabaikan terhadap ukuran Lebesgue,
dan bahwa fungsi Cantor\cmmnt{ (134H)} merupakan fungsi
tak menurun
kontinu $f:[0,1]\to[0,1]$ sedemikian sehingga $f(0)=0$, $f(1)=1$ dan
$f$ konstan pada setiap interval yang menyusun $[0,1]\setminus C$.
Akibatnya, jika kita menetapkan $g(x)=\bover12(x+f(x))$ untuk $x\in[0,1]$, maka
$g:[0,1]\to[0,1]$ merupakan permutasi kontinu\cmmnt{ sedemikian sehingga ukuran
Lebesgue dari $g[C]$ adalah $\bover12$ (134I);
akibatnya $g^{-1}:[0,1]\to[0,1]$
kontinu}.   Sekarang perluas $g$ menjadi permutasi $h:\Bbb R\to\Bbb R$ dengan
menetapkan $h(x)=x$ untuk $x\in\Bbb R\setminus[0,1]$.
\cmmnt{Maka $h$ dan $h^{-1}$ kontinu.   Perhatikan bahwa $h[C]=g[C]$
mempunyai ukuran Lebesgue $\bover12$.}

Misalkan $\nu_1$ ukuran integral tak tentu yang didefinisikan dari ukuran
Lebesgue $\mu$ pada $\Bbb R$ dan
fungsi $2\chi(h[C])$\cmmnt{;  yaitu, $\nu_1E=2\mu(E\cap h[C])$ kapan pun
ungkapan ini terdefinisi}.   \cmmnt{Menurut 256E,} $\nu_1$ merupakan ukuran Radon, dan
$\nu_1h[C]=\nu_1\Bbb R=1$.   Misalkan $\nu$ ukuran $\nu_1(h^{-1})^{-1}$,
yaitu,
$\nu E=\nu_1h[E]$ tepat bagi himpunan-himpunan $E\subseteq\Bbb R$ sedemikian sehingga
$h[E]\in\dom\nu_1$.   Maka $\nu$ merupakan ukuran probabilitas Radon pada
$\Bbb R$\cmmnt{, menurut 256G}, dan $\nu C=1$,
$\nu(\Bbb R\setminus C)=\mu C=0$.

\cmmnt{
\leader{256I}{Catatan (a)} Ukuran $\nu$ dalam 256Hc, yang kadang-kadang disebut
{\bf ukuran Cantor}, merupakan contoh klasik, dan karena itu mempunyai banyak
konstruksi, beberapa di antaranya jauh lebih alami daripada yang saya gunakan di sini (lihat
256Xk, dan juga 264Ym di bawah).   Tetapi saya memilih metode di atas
karena metode itu
secara langsung, tanpa penyelidikan lebih lanjut atau mengacu pada teori umum yang lebih
lanjut, menghasilkan fakta bahwa $\nu$ merupakan ukuran Radon.

\spheader 256Ib Contoh-contoh di atas dipilih untuk mewakili
kasus-kasus ekstrem dalam `dekomposisi Lebesgue' yang diuraikan dalam
232I.   Jika $\nu$ merupakan ukuran Radon (berhingga total) pada $\BbbR^r$, kita
dapat menggunakan 232Ib untuk menyatakan pembatasannya $\nu\restr\Cal B$ pada
aljabar-sigma Borel sebagai $\nu_p+\nu_{ac}+\nu_{cs}$, dengan $\nu_p$ sebagai
bagian `massa-titik' atau `atomik' dari $\nu\restr\Cal B$, $\nu_{ac}$ sebagai
bagian `kontinu
mutlak' (terhadap ukuran Lebesgue), dan $\nu_{cs}$ sebagai
`bagian singular tanpa atom'.   Dalam contoh 256Hb, kita mempunyai
$\nu\restr\Cal B=\nu_p$;  dalam 256E, jika kita mulai dari ukuran Lebesgue, kita
mempunyai $\nu\restr\Cal B=\nu_{ac}$;
dan dalam 256Hc kita mempunyai $\nu\restr\Cal B=\nu_{cs}$.
}%end of comment

\leader{256J}{Ukuran Radon \dvrocolon{yang kontinu mutlak}}\cmmnt{
Ada baiknya berhenti sejenak untuk menelaah ukuran-ukuran integral tak tentu yang
diuraikan dalam 256E.

\medskip

\noindent}{\bf Proposisi} Misalkan $\nu$ suatu ukuran Radon
pada $\BbbR^r$, dan tuliskan $\mu$ untuk ukuran Lebesgue pada
$\BbbR^r$.   Maka pernyataan-pernyataan berikut ekuivalen:

\quad(i) $\nu$ merupakan ukuran integral tak tentu terhadap $\mu$;

\quad(ii) $\nu E=0$ kapan pun $E$ merupakan subhimpunan Borel dari $\BbbR^r$ dan
$\mu E=0$.

\noindent Dalam hal ini, jika $g\in\eusm L^0(\mu)$ dan
$\int_Eg\,d\mu=\nu E$ untuk
setiap himpunan Borel $E\subseteq\BbbR^r$, maka $g$ merupakan turunan Radon-Nikod\'ym
dari $\nu$ terhadap $\mu$\cmmnt{ dalam arti
232Hf}.

\proof{{\bf (a)(i)$\Rightarrow$(ii)} Jika $f$ merupakan turunan Radon-Nikod\'ym
dari $\nu$ terhadap $\mu$, maka tentu saja

\Centerline{$\nu E=\int_Efd\mu=0$}

\noindent kapan pun $\mu E=0$.

\medskip

\quad{\bf (ii)$\Rightarrow$(i)} Jika $\nu E=0$ untuk setiap himpunan Borel $E$ yang
terabaikan terhadap $\mu$, maka $\nu E$ terdefinisi dan sama dengan $0$ untuk
setiap himpunan $E$ yang terabaikan terhadap $\mu$, karena $\nu$ lengkap dan setiap
himpunan yang terabaikan terhadap $\mu$ termuat dalam suatu himpunan Borel yang terabaikan terhadap $\mu$.
Akibatnya $\dom\nu$ memuat
domain $\Sigma$ dari $\mu$, karena setiap himpunan terukur Lebesgue dapat
dinyatakan sebagai gabungan suatu himpunan Borel dan suatu himpunan terabaikan.

Untuk setiap $n\in\Bbb N$ tetapkan $E_n=\{x:n\le\|x\|<n+1\}$, sehingga
$\sequencen{E_n}$ merupakan partisi $\BbbR^r$ menjadi himpunan-himpunan Borel terbatas.
Tetapkan $\nu_nE=\nu(E\cap E_n)$ untuk setiap himpunan terukur Lebesgue $E$ dan
setiap $n\in\Bbb N$.   Sekarang $\nu_n$ kontinu mutlak terhadap
$\mu$ (232Ba), sehingga menurut teorema Radon-Nikod\'ym (232F)
terdapat fungsi yang dapat diintegralkan terhadap $\mu$, yaitu $f_n$, sedemikian sehingga
$\int_Ef_nd\mu=\nu_n E$
untuk setiap himpunan terukur Lebesgue $E$.   Karena $\nu_n E\ge 0$ untuk
setiap $E\in\Sigma$,
$f_n\ge 0$ hampir di mana-mana;  karena $\nu_n(\BbbR^r\setminus E_n)=0$, $f_n=0$
hampir di mana-mana\ pada $\BbbR^r\setminus E_n$.   Sekarang, jika kita menetapkan

\Centerline{$f=\max(0,\sum_{n=0}^{\infty}f_n)$,}

\noindent $f$ akan terdefinisi hampir di mana-mana terhadap $\mu$ dan kita akan mempunyai

\Centerline{$\int_Efd\mu
=\sum_{n=0}^{\infty}\int_Ef_nd\mu
=\sum_{n=0}^{\infty}\nu(E\cap E_n)=\nu E$}

\noindent untuk setiap himpunan Borel $E$, sehingga ukuran integral tak tentu
$\nuprime$ yang didefinisikan oleh $f$ dan $\mu$ berimpit dengan $\nu$ pada
himpunan-himpunan Borel.   Karena ini memastikan bahwa $\nuprime$ berhingga secara lokal,
$\nuprime$ merupakan ukuran Radon, menurut 256E, dan sama dengan $\nu$, menurut 256D.
Dengan demikian $\nu$ merupakan ukuran integral tak tentu terhadap $\mu$.

\medskip

{\bf (b)} Seperti pada (a-ii) di atas, $g$ harus terintegralkan secara lokal dan ukuran
integral tak tentu yang didefinisikan oleh $g$ berimpit dengan $\nu$ pada
himpunan-himpunan Borel, sehingga identik dengan $\nu$.
}%end of proof of 256J

\leader{256K}{\dvrocolon{Produk}}\cmmnt{ Kelas ukuran Radon
pada ruang Euklides stabil terhadap beragam operasi, sebagaimana telah
kita lihat;  khususnya, kita memperoleh hasil berikut.

\medskip

\noindent}{\bf Teorema} Misalkan $\nu_1$, $\nu_2$ masing-masing ukuran Radon
pada $\BbbR^r$ dan $\BbbR^s$.
Misalkan $\lambda$ ukuran produk lengkap yang ditentukan secara lokal (c.l.d.)\ dari keduanya pada
$\BbbR^r\times\BbbR^s$.   Maka $\lambda$ merupakan ukuran Radon.

\cmmnt{\medskip

\noindent{\bf Catatan} Ketika saya mengatakan bahwa $\lambda$ bersifat `Radon' menurut
definisi dalam 256A, tentu saja saya mengidentifikasi
$\BbbR^r\times\BbbR^s$ dengan $\BbbR^{r+s}$, seperti dalam 251M-251N.
}

\proof{ Saya berharap notasi berikut terasa alami.
Tuliskan $\Sigma_1$, $\Sigma_2$
untuk domain $\nu_1$, $\nu_2$;  $\Cal B_r$, $\Cal B_s$ untuk
aljabar-sigma Borel dari $\BbbR^r$, $\BbbR^s$;  $\Lambda$ untuk
domain $\lambda$;  dan $\Cal B$ untuk aljabar-sigma Borel dari
$\BbbR^{r+s}$.

Karena setiap $\nu_i$ merupakan pelengkapan dari pembatasannya pada himpunan-himpunan Borel
(256C), $\lambda$ merupakan produk dari $\nu_1\restr\Cal B_r$ dan
$\nu_2\restr\Cal B_s$ (251T).   Karena $\nu_1\restr\Cal B_r$ dan
$\nu_2\restr\Cal B_s$ bersifat $\sigma$-hingga (256Ba, 212Ga), $\lambda$ harus
merupakan pelengkapan dari pembatasannya pada $\Cal B_r\tensorhat\Cal B_s$,
yang menurut 251M diidentifikasi dengan $\Cal B$.   Dengan menetapkan
$Q_n=\{(x,y):\|x\|\le n$, $\|y\|\le n\}$ kita mempunyai

\Centerline{$\lambda Q_n=\nu_1\{x:\|x\|\le n\}\cdot\nu_2\{y:\|y\|\le n\}
<\infty$}

\noindent untuk setiap $n$, sedangkan setiap subhimpunan terbatas dari $\BbbR^{r+s}$
termuat dalam suatu $Q_n$.   Jadi $\lambda\restr\Cal B$ berhingga secara lokal,
dan pelengkapannya $\lambda$ merupakan ukuran Radon, menurut 256C.
}%end of proof of 256K

\cmmnt{
\leader{256L}{Catatan} Kita melihat dari 253I bahwa jika $\nu_1$ dan
$\nu_2$ masing-masing merupakan ukuran Radon pada $\BbbR^r$ dan $\BbbR^s$, dan
keduanya merupakan ukuran integral tak tentu terhadap ukuran Lebesgue, maka
ukuran produk keduanya pada $\BbbR^{r+s}$ juga merupakan ukuran integral tak tentu
terhadap ukuran Lebesgue.
}%end of comment

\leader{*256M}{}\cmmnt{ Demi penerapan dalam \S286 di bawah, saya
menyertakan satu hasil lagi, yang sebenarnya merupakan salah satu sifat
mendasar ukuran Radon, sebagaimana akan tampak dalam \S414.

\medskip

\noindent}{\bf Proposisi} Misalkan $\nu$ ukuran Radon pada $\Bbb R^r$,
dan $D$ sebarang subhimpunan $\Bbb R^r$. Misalkan $\Phi$ keluarga tak kosong
yang terarah ke atas dari fungsi-fungsi kontinu nonnegatif dari $D$ ke
$\Bbb R$. Untuk $x\in D$ tetapkan $g(x)=\sup_{f\in\Phi}f(x)$ dalam
$[0,\infty]$. Maka

(a) $g:D\to[0,\infty]$ semikontinu bawah,
dan karena itu terukur Borel;

(b) $\int_Dg\,d\nu=\sup_{f\in\Phi}\int_Dfd\nu$.

\proof{{\bf (a)} Untuk setiap $u\in[-\infty,\infty]$,

\Centerline{$\{x:x\in D,\,g(x)>u\}
=\bigcup_{f\in\Phi}\{x:x\in D,\,f(x)>u\}$}

\noindent merupakan himpunan terbuka untuk topologi subruang pada $D$
(2A3C), sehingga merupakan irisan $D$ dengan suatu subhimpunan Borel dari
$\Bbb R^r$. Ini cukup untuk menunjukkan bahwa $g$ terukur Borel
(121B-121C).

\medskip

{\bf (b)} Karena itu $\int_Dg\,d\nu$ terdefinisi dalam $[0,\infty]$,
dan tentu saja $\int_Dg\,d\nu\ge\sup_{f\in\Phi}\int_Dfd\nu$.

Untuk ketaksamaan sebaliknya, perhatikan bahwa terdapat himpunan terhitung
$\Psi\subseteq\Phi$ sedemikian sehingga $g(x)=\sup_{f\in\Psi}f(x)$ untuk
setiap $x\in D$. \Prf\ Untuk $a\in\Bbb Q$, $q$, $q'\in\BbbQ^r$ tetapkan

\Centerline{$\Phi_{aqq'}
=\{f:f\in\Phi,\,f(y)>a\text{ kapan pun }y\in D\cap[q,q']\}$,}

\noindent dengan menafsirkan $[q,q']$ seperti dalam 115G. Pilih
$f_{aqq'}\in\Phi_{aqq'}$ jika $\Phi_{aqq'}$ tidak kosong, dan secara sebarang
dalam $\Phi$ jika tidak; serta tetapkan
$\Psi=\{f_{aqq'}:a\in\Bbb Q,\,q,\,q'\in\BbbQ^r\}$, sehingga $\Psi$ merupakan
subhimpunan terhitung dari $\Phi$. Jika $x\in D$ dan $b<g(x)$, terdapat
$a\in\Bbb Q$ sedemikian sehingga $b\le a<g(x)$; terdapat
$\hat f\in\Phi$ sedemikian sehingga $\hat f(x)>a$; karena $\hat f$
kontinu, terdapat $q$, $q'\in\Bbb Q^r$ sedemikian sehingga $q\le x\le q'$
dan $\hat f(y)>a$ kapan pun $y\in D\cap[q,q']$; sehingga
$\hat f\in\Phi_{aqq'}$, $\Phi_{aqq'}\ne\emptyset$,
$f_{aqq'}\in\Phi_{aqq'}$ dan $\sup_{f\in\Psi}f(x)\ge f_{aqq'}(x)\ge b$.
Karena $b$
sebarang, $g(x)=\sup_{f\in\Psi}f(x)$.\ \Qed

Misalkan $\sequencen{f_n}$ suatu barisan yang mencakup $\Psi$. Karena $\Phi$
terarah ke atas, kita dapat memilih $\sequencen{f'_n}$ dalam $\Phi$ secara
induktif sedemikian sehingga $f'_{n+1}\ge\max(f'_n,f_n)$ untuk setiap
$n\in\Bbb N$. Jadi $\sequencen{f'_n}$ merupakan barisan tak menurun dalam
$\Phi$ dan
$\sup_{n\in\Bbb N}f'_n(x)\ge\sup_{f\in\Psi}f(x)=g(x)$ untuk setiap
$x\in D$. Menurut teorema B.Levi,

\Centerline{$\int_Dg\,d\nu
\le\sup_{n\in\Bbb N}\int_Df'_nd\nu\le\sup_{f\in\Phi}\int_Dfd\nu$,}

\noindent dan kita memperoleh ketaksamaan yang diperlukan.
}%end of proof of 256M

\exercises{
\leader{256X}{Latihan dasar $\pmb{>}$(a)}
%\spheader 256Xa
Misalkan $\nu$ ukuran pada $\BbbR^r$. (i) Tunjukkan bahwa ukuran itu
berhingga secara lokal, dalam arti 256Ab, jika dan hanya jika untuk setiap
$x\in\BbbR^r$ terdapat $\delta>0$ sedemikian sehingga
$\nu^*B(x,\delta)<\infty$. \Hint{himpunan-himpunan
$B(\tbf{0},n)$ kompak.} (ii) Tunjukkan bahwa dalam hal ini $\nu$
$\sigma$-hingga.
%256A

\sqheader 256Xb Misalkan $\nu$ ukuran Radon pada $\BbbR^r$ dan $\Cal G$
keluarga tak kosong
yang terarah ke atas dari himpunan-himpunan terbuka dalam $\BbbR^r$.
(i) Tunjukkan bahwa
$\nu(\bigcup\Cal G)=\sup_{G\in\Cal G}\nu G$. \Hint{perhatikan bahwa jika
$K\subseteq\bigcup\Cal G$ kompak, maka $K\subseteq G$ untuk suatu
$G\in\Cal G$.} (ii) Tunjukkan bahwa
$\nu(E\cap\bigcup\Cal G)=\sup_{G\in\Cal G}\nu(E\cap G)$ untuk setiap
himpunan $E$ yang terukur terhadap $\nu$.
%256A

\sqheader 256Xc Misalkan $\nu$ ukuran Radon pada $\BbbR^r$ dan $\Cal F$
keluarga tak kosong yang terarah ke bawah dari himpunan-himpunan tertutup dalam
$\BbbR^r$ sedemikian sehingga $\inf_{F\in\Cal F}\nu F<\infty$.
(i) Tunjukkan bahwa
$\nu(\bigcap\Cal F)=\inf_{F\in\Cal F}\nu F$. \Hint{terapkan 256Xb(ii) pada
$\Cal G=\{\BbbR^r\setminus F:F\in\Cal F\}$.} (ii) Tunjukkan bahwa
$\nu(E\cap\bigcap\Cal F)=\inf_{F\in\Cal F}\nu(E\cap F)$ untuk setiap
$E$ dalam domain $\nu$.
%256A

\sqheader 256Xd Tunjukkan bahwa ukuran Radon $\nu$ pada $\BbbR^r$ tanpa atom
jika dan hanya jika $\nu\{x\}=0$ untuk setiap $x\in\BbbR^r$.
\Hint{terapkan 256Xc dengan
$\Cal F=\{F:F\subseteq E$ tertutup, tidak terabaikan$\}$.}
%256A

\spheader 256Xe Misalkan $\nu_1$, $\nu_2$ ukuran Radon pada
$\BbbR^r$, dan $\alpha_1$, $\alpha_2\in\ooint{0,\infty}$. Tetapkan
$\Sigma=\dom\nu_1\cap\dom\nu_2$, dan untuk $E\in\Sigma$ tetapkan
$\nu E=\alpha_1\nu_1E+\alpha_2\nu_2E$. Tunjukkan bahwa $\nu$ merupakan
ukuran Radon pada $\BbbR^r$. Tunjukkan bahwa $\nu$ merupakan ukuran
integral tak tentu terhadap ukuran Lebesgue jika dan hanya jika $\nu_1$, $\nu_2$
demikian, dan bahwa dalam hal ini suatu kombinasi linear
dari turunan-turunan Radon-Nikod\'ym dari $\nu_1$ dan $\nu_2$ merupakan
turunan Radon-Nikod\'ym dari $\nu$.
%256A

\sqheader 256Xf Misalkan $\nu$ ukuran Radon pada $\BbbR^r$. (i) Tunjukkan
bahwa terdapat tepat satu himpunan tertutup $F\subseteq\BbbR^r$ sedemikian
sehingga, untuk himpunan terbuka $G\subseteq\BbbR^r$, $\nu G>0$ jika dan
hanya jika $G\cap F\ne\emptyset$.
($F$ disebut {\bf dukungan} dari $\nu$.) (ii) Secara umum, suatu
himpunan $A\subseteq\BbbR^r$ disebut {\bf menopang dirinya sendiri} jika
$\nu^*(A\cap G)>0$ kapan pun $G\subseteq\BbbR^r$ merupakan himpunan terbuka
yang bertemu dengan $A$. Tunjukkan bahwa untuk setiap himpunan tertutup
$F\subseteq\BbbR^r$ terdapat tepat satu himpunan tertutup yang menopang dirinya sendiri
$F'\subseteq F$ sedemikian sehingga $\nu(F\setminus F')=0$.
%256A

\sqheader 256Xg Tunjukkan bahwa ukuran $\nu$ pada $\Bbb R$ merupakan
ukuran Radon jika dan hanya jika ukuran itu merupakan ukuran Lebesgue--Stieltjes
sebagaimana diuraikan dalam 114Xa. Tunjukkan bahwa dalam hal ini $\nu$
merupakan ukuran integral tak tentu terhadap ukuran Lebesgue jika dan hanya jika
fungsi $x\mapsto\nu[a,x]:[a,b]\to\Bbb R$
kontinu mutlak kapan pun $a\le b$ dalam $\Bbb R$.
%256C

\spheader 256Xh Misalkan $\nu$ ukuran Radon pada $\BbbR^r$. Misalkan
$C_k$ ruang fungsi kontinu bernilai real pada
$\BbbR^r$ dengan dukungan terbatas. Tunjukkan bahwa untuk setiap
fungsi yang dapat diintegralkan terhadap $\nu$, sebutlah $f$, dan setiap $\epsilon>0$
terdapat $g\in C_k$ sedemikian sehingga $\int|f-g|d\nu\le\epsilon$.
\Hint{gunakan argumen dari 242O, tetapi dalam (a-i) dari buktinya mulailah
dengan interval-interval {\it tertutup} $I$.}
%256C

\spheader 256Xi Misalkan $\nu$ ukuran Radon pada $\BbbR^r$, dan $\nu^*$
ukuran luar yang bersesuaian. Tunjukkan bahwa
$\nu^*A=\inf\{\nu G:G\supseteq A$ terbuka$\}$ untuk setiap himpunan
$A\subseteq\BbbR^r$.
%256D

\spheader 256Xj Misalkan $\nu$, $\nuprime$ dua ukuran Radon pada
$\BbbR^r$, dan
andaikan bahwa $\nu I=\nuprime I$ untuk setiap interval setengah terbuka
$I\subseteq\BbbR^r$ (definisi: 115Ab). Tunjukkan bahwa $\nu=\nuprime$.
%256D

\spheader 256Xk Misalkan $\nu$ ukuran Cantor (256Hc).
(i) Tunjukkan bahwa jika $C_n$ adalah himpunan ke-$n$ yang digunakan dalam
konstruksi himpunan Cantor, sehingga $C_n$ terdiri atas $2^n$ interval dengan
panjang $3^{-n}$, maka $\nu I=2^{-n}$ untuk setiap interval $I$ yang menyusun
$C_n$. (ii) Misalkan $\lambda$ ukuran biasa pada $\{0,1\}^{\Bbb N}$
(254J). Definisikan $\phi:\{0,1\}^{\Bbb N}\to \Bbb R$ dengan menetapkan
$\phi(x)=\bover23\sum_{n=0}^{\infty}3^{-n}x(n)$ untuk setiap
$x\in\{0,1\}^{\Bbb N}$. Tunjukkan bahwa $\phi$ merupakan bijeksi antara
$\{0,1\}^{\Bbb N}$ dan $C$. (iii) Tunjukkan bahwa jika $\Cal B$ merupakan
aljabar-sigma Borel dari $\Bbb R$, maka $\{\phi^{-1}[E]:E\in\Cal B\}$
tepat merupakan aljabar-sigma subhimpunan $\{0,1\}^{\Bbb N}$
yang dibangkitkan oleh himpunan-himpunan $\{x:x(n)=i\}$ untuk
$n\in\Bbb N$, $i\in\{0,1\}$. (iv) Tunjukkan bahwa $\phi$ merupakan
isomorfisme antara $(\{0,1\}^{\Bbb N},\lambda)$ dan $(C,\nu_C)$, dengan
$\nu_C$ ukuran subruang pada $C$ yang diinduksi oleh $\nu$.
%256H

\spheader 256Xl Misalkan $\nu$ dan $\nuprime$ dua ukuran Radon pada
$\BbbR^r$. Tunjukkan bahwa $\nuprime$ merupakan ukuran integral tak tentu terhadap
$\nu$ jika dan hanya jika $\nuprime E=0$ kapan pun $\nu E=0$, dan dalam hal
ini fungsi $f$ merupakan turunan Radon-Nikod\'ym dari $\nuprime$ terhadap
$\nu$ jika dan hanya jika $\int_Efd\nu=\nuprime E$ untuk setiap himpunan
Borel $E$.
%256J

\leader{256Y}{Latihan lanjutan (a)}
%\spheader 256Ya
Misalkan $\nu$ ukuran Radon pada $\BbbR^r$, dan $X$ sebarang subhimpunan
$\Bbb R^r$; misalkan $\nu_X$ ukuran subruang pada $X$ dan $\Sigma_X$
domainnya, dan lengkapi $X$ dengan
topologi subruangnya. Tunjukkan bahwa $\nu_X$ mempunyai sifat-sifat berikut:
(i) $\nu_X$ lengkap dan ditentukan secara lokal; (ii) setiap
subhimpunan terbuka dari $X$ termasuk dalam $\Sigma_X$; (iii)
$\nu_XE=\sup\{\nu_XF:F\subseteq E$ tertutup dalam $X\}$ untuk setiap
$E\in\Sigma_X$; (iv) kapan pun $\Cal G$ merupakan keluarga tak kosong
subhimpunan terbuka dari $X$ yang terarah ke atas,
$\nu_X(\bigcup\Cal G)=\sup_{G\in\Cal G}\nu_XG$;
(v) setiap titik di $X$ termasuk dalam suatu himpunan terbuka yang ukurannya berhingga.
%256A

\spheader 256Yb Misalkan $\nu$ ukuran Radon pada $\BbbR^r$, dengan domain
$\Sigma$, dan $f:\Bbb R^r\to\Bbb R$ suatu fungsi. Tunjukkan bahwa pernyataan
berikut ekuivalen: (i) $f$ terukur terhadap $\Sigma$;
(ii) untuk setiap himpunan tak terabaikan
$E\in\Sigma$ terdapat himpunan tak terabaikan $F\in\Sigma$ sedemikian
sehingga $F\subseteq E$ dan $f\restr F$ kontinu; (iii) untuk setiap himpunan
$E\in\Sigma$, $\nu E=\sup_{K\in\Cal K_f,K\subseteq E}\nu K$, dengan
$\Cal K_f=\{K:K\subseteq\BbbR^r$ kompak, $f\restr K$ kontinu$\}$.
\Hint{untuk (ii)$\Rightarrow$(i), terapkan 215B(iv) pada $\Cal K_f$.}
%256F

\spheader 256Yc Ambil $\nu$, $X$, $\nu_X$ dan $\Sigma_X$ seperti dalam 256Ya.
Andaikan bahwa
$f:X\to\Bbb R$ suatu fungsi. Tunjukkan bahwa $f$ terukur terhadap $\Sigma_X$
jika dan hanya jika untuk setiap himpunan terukur tak terabaikan
$E\subseteq X$ terdapat himpunan terukur tak terabaikan $F\subseteq E$
sedemikian sehingga $f\restr F$ kontinu.
%256F

\spheader 256Yd\dvAnew{2011}(i) Misalkan $\lambda$ ukuran biasa pada
$\{0,1\}^{\Bbb N}$. Definisikan $\psi:\{0,1\}^{\Bbb N}\to\{0,1\}^{\Bbb N}$
dengan menetapkan $\psi(x)(i)=x(i+1)$ untuk $x\in\{0,1\}^{\Bbb N}$ dan
$i\in\Bbb N$. Tunjukkan bahwa $\psi$ bersifat \imp. (ii) Definisikan
$\theta:\Bbb R\to\Bbb R$ dengan menetapkan
$\theta(t)=\fraction{3t}=3t-\lfloor 3t\rfloor$ untuk $t\in\Bbb R$.
Tunjukkan bahwa $\theta$ bersifat \imp\ untuk ukuran Cantor sebagaimana
didefinisikan dalam 256Hc.
%256H 256Xk

\spheader 256Ye Misalkan $\sequencen{\nu_n}$ barisan ukuran Radon
pada $\BbbR^r$. Tunjukkan bahwa terdapat ukuran Radon $\nu$ pada $\BbbR^r$
sedemikian sehingga setiap $\nu_n$ merupakan ukuran integral tak tentu terhadap
$\nu$. \Hint{carilah barisan $\sequencen{\alpha_n}$ dari
bilangan-bilangan positif ketat sedemikian sehingga
$\sum_{n=0}^{\infty}\alpha_n\nu_n
B(\tbf{0},k)<\infty$ untuk setiap $k$, dan tetapkan
$\nu=\sum_{n=0}^{\infty}\alpha_n\nu_n$, dengan menggunakan gagasan 256Xe.}
%256J, 256Xe, 256Xl

\spheader 256Yf Suatu himpunan
$G\subseteq\BbbR^{\Bbb N}$ disebut {\bf terbuka} jika untuk setiap $x\in G$
terdapat $n\in\Bbb N$, $\delta>0$ sedemikian sehingga

\Centerline{$\{y:y\in\BbbR^{\Bbb N},\,|y(i)-x(i)|<\delta$ untuk setiap
$i\le n\}\subseteq G$.}

\noindent {\bf Aljabar-sigma Borel} dari $\BbbR^{\Bbb N}$ adalah
aljabar-sigma $\Cal B$ dari subhimpunan-subhimpunan
$\BbbR^{\Bbb N}$ yang dibangkitkan, dalam arti 111Gb, oleh keluarga
$\frak T$ dari himpunan-himpunan terbuka. (i) Tunjukkan bahwa $\frak T$
merupakan topologi (2A3A).
(ii) Tunjukkan bahwa filter $\Cal F$ pada $\BbbR^{\Bbb N}$ konvergen ke
$x\in\BbbR^{\Bbb N}$ jika dan hanya jika $\pi_i[[\Cal F]]\to x(i)$ untuk
setiap $i\in\Bbb N$, dengan
$\pi_i(y)=y(i)$ untuk $i\in\Bbb N$, $y\in\BbbR^{\Bbb N}$. (iii)
Tunjukkan bahwa $\Cal B$ merupakan aljabar-sigma yang dibangkitkan oleh
himpunan-himpunan berbentuk
$\{x:x\in \BbbR^{\Bbb N},\,x(i)\le a\}$,
dengan $i$ menjangkau $\Bbb N$ dan $a$ menjangkau $\Bbb R$.
(iv) Tunjukkan bahwa jika $\alpha_i\ge 0$ untuk setiap $i\in\Bbb N$, maka
$\{x:|x(i)|\le\alpha_i\Forall i\in\Bbb N\}$ kompak.
%2A3R
(v) Tunjukkan bahwa setiap himpunan terbuka dalam $\BbbR^{\Bbb N}$ merupakan
gabungan suatu barisan himpunan tertutup. \Hint{perhatikan himpunan-himpunan
berbentuk
$\{x:q_i\le x(i)\le q'_i\Forall i\le n\}$, dengan $q_i$, $q'_i\in\Bbb Q$
untuk $i\le n$.} (vi) Tunjukkan bahwa jika $\nu_0$ sebarang ukuran
probabilitas dengan domain $\Cal B$, maka pelengkapannya, $\nu$, regular dalam
terhadap himpunan-himpunan kompak, dan karena itu dapat disebut `ukuran Radon
pada $\BbbR^{\Bbb N}$'. \Hint{tunjukkan bahwa terdapat himpunan-himpunan
kompak yang ukurannya sedekat yang diinginkan dengan $1$, dan karena itu setiap
himpunan terbuka, dan setiap himpunan tertutup, memuat suatu himpunan
K$_{\sigma}$ dengan ukuran yang sama.}
%256K  256Cprf
}%end of exercises


\endnotes{
\Notesheader{256} Ukuran Radon pada ruang Euklides sangat khusus,
dan hasil-hasil bagian ini tidak memberikan petunjuk yang jelas mengenai
arah perkembangan teori ketika diterapkan pada jenis ruang topologis lain.
Dengan bahan di sini Anda dapat mencoba mengembangkan teori ukuran Radon pada
ruang metrik lengkap separabel, asalkan Anda menggunakan 256Xa sebagai dasar
definisi `berhingga secara lokal'. Pada ruang-ruang inilah suatu versi 256C berlaku.
(Lihat 256Yf.) Namun, untuk perumuman ke
jenis ruang topologis lain, dan untuk bagian-bagian teori yang lebih menarik
pada $\BbbR^r$, saya harus meminta Anda menunggu
hingga Jilid 4.
Tujuan saya memperkenalkan ukuran Radon di sini sangat
terbatas; saya hanya ingin memberikan landasan bagi \S257 dan \S271 yang cukup
kokoh sehingga tidak
memerlukan revisi kemudian. Sebenarnya saya kira yang sungguh-sungguh kita
perlukan hanyalah ukuran probabilitas Radon.

Kesulitan teknis utama dalam definisi
`ukuran Radon' di sini terletak
pada tuntutan kelengkapan. Sangat mungkin bahwa untuk
segala hal yang dikaji dalam jilid ini, akan lebih sederhana jika kita melihat
ukuran berhingga secara lokal dengan domain aljabar himpunan-himpunan Borel.
Hal ini akan memaksa kita menggunakan sejumlah uraian tak langsung ketika
menangani ukuran Lebesgue itu sendiri dan turunannya, karena ukuran Lebesgue
didefinisikan pada aljabar-sigma yang lebih besar; tetapi keberatan yang serius
muncul dalam teori yang lebih lanjut, ketika himpunan-himpunan non-Borel dari
berbagai jenis menjadi pusat perhatian. Karena tujuan saya dalam buku ini
adalah memberikan landasan yang kokoh bagi kajian semua segi teori ukuran, saya
meminta Anda bersusah payah sedikit lebih banyak sekarang agar terhindar dari
kemungkinan harus menyusun kembali seluruh gagasan Anda nanti.
Kesulitan tambahan ini muncul, misalnya, dalam 256D, 256Xe dan 256Xj; karena
ukuran Radon yang berbeda didefinisikan pada aljabar-sigma yang berbeda, kita
harus memeriksa bahwa dua ukuran Radon yang bersepakat pada himpunan-himpunan
kompak, atau pada himpunan-himpunan terbuka, memiliki domain yang sama. Sebagai
keuntungannya, sebagian kekuatan 256G timbul dari kenyataan bahwa ukuran citra
Radon $\nu\phi^{-1}$ didefinisikan pada seluruh aljabar-sigma
$\{F:\phi^{-1}[F]\in\dom(\nu)\}$, bukan hanya pada himpunan Borel.

Butir teknis selanjutnya, bahwa ukuran Radon diharapkan berhingga secara lokal,
menimbulkan lebih sedikit kesulitan; akibatnya ialah bahwa dari kebanyakan
sudut pandang hanya terdapat sedikit perbedaan antara ukuran Radon umum dan
ukuran Radon berhingga total. Syarat tambahan yang dengan jelas harus
dimasukkan ke dalam hipotesis hasil-hasil seperti 256E dan
256G tidak membebani intuisi maupun ingatan.

Pada dasarnya, kita mempunyai dua definisi ukuran Radon pada ruang Euklides:
ukuran Radon adalah ukuran topologis berhingga secara lokal yang regular dalam,
dan juga merupakan pelengkapan ukuran Borel berhingga secara lokal. Ekuivalensi kedua
definisi ini adalah Teorema 256C. Definisi yang terakhir lebih sesuai untuk
256K, dan yang pertama untuk 256G. `Regularitas dalam' dalam definisi dasar
merujuk pada himpunan kompak; kita juga mempunyai bentuk regularitas dalam
terhadap himpunan tertutup (256Bb) dan himpunan K$_{\sigma}$ (256Bc), serta
gagasan pelengkap berupa `regularitas luar' terhadap himpunan terbuka (256Xi).

}%end of notes

\discrpage
